\section{Diseño e implementación de la API}
\label{sec:api}

\subsection{Contrato}

Todos los recursos viven bajo el prefijo \codigo{/api/v1}. La especificación
OpenAPI la genera FastAPI a partir de los propios tipos y está versionada en
\codigo{docs/openapi.json}; el CI falla si diverge del código, de modo que la
documentación no puede desincronizarse en silencio. \estado{verificado}

\begin{table}[htbp]
\centering
\caption{Endpoints implementados.}
\label{tab:endpoints}
\scriptsize
\begin{tabularx}{\textwidth}{@{}llcXl@{}}
\toprule
\textbf{Método} & \textbf{Ruta} & \textbf{Auth} & \textbf{Errores previstos} &
\textbf{Éxito} \\
\midrule
POST & \codigo{/auth/register} & no &
  409 correo repetido, 413, 422, 429 & 201 \\
POST & \codigo{/auth/login} & no & 401, 413, 429 & 200 \\
GET & \codigo{/accounts/me} & sí & 401 & 200 \\
GET & \codigo{/accounts/me/balance} & sí & 401 & 200 \\
GET & \codigo{/accounts/me/transactions} & sí & 401, 422 & 200 \\
POST & \codigo{/transfers} & sí &
  400 falta la clave, 401, 404 destino inexistente, 409 clave en curso, 413,
  422 fondos o clave reutilizada, 429, 503 & 201 \\
GET & \codigo{/transfers/\{id\}} & sí & 401, 404 & 200 \\
GET & \codigo{/health} & no & --- & 200 \\
\bottomrule
\end{tabularx}
\end{table}

\subsection{Tres decisiones del contrato}

\paragraph{Los importes nunca viajan como decimal.} El campo
\codigo{amount\_cents} es un entero de céntimos. Para la interfaz se añade
\codigo{amount\_display}, una cadena ya formateada. Si llega un decimal, la
petición se rechaza con 422 en lugar de redondearse: redondear en silencio
sería decidir por el usuario sobre su dinero. \estado{verificado}

\paragraph{404 en lugar de 403 para recursos ajenos.} Cuando alguien pide una
operación que existe pero en la que no participa, la respuesta es 404. Un 403
confirmaría la existencia del recurso y permitiría enumerar operaciones
ajenas, que es la vía de explotación típica de la autorización rota a nivel
de objeto~\cite{owaspapi2023}.

\paragraph{Forma uniforme del error.} Todas las respuestas de fallo comparten
estructura y ninguna incluye traza interna. El \codigo{request\_id} aparece
también en el registro estructurado, lo que permite correlacionar sin exponer
nada.

\begin{lstlisting}[language=,caption={Cuerpo de error, idéntico para todos los
fallos.},label={lst:error}]
{
  "error": {
    "code": "insufficient_funds",
    "message": "El saldo disponible no cubre el importe solicitado.",
    "request_id": "85006c24b5494c3c89ecdc64df8a5cb6",
    "details": { "transfer_id": "...", "balance_cents": 7000 }
  }
}
\end{lstlisting}

\subsection{Observabilidad}

Registro estructurado en JSON con un identificador de correlación por
petición, propagado por un middleware y devuelto en la cabecera
\codigo{X-Request-ID}. Hay una lista de denegación explícita: nunca se
registran contraseñas, tokens, cabeceras de autorización ni claves de
idempotencia completas. Los importes y los identificadores de cuenta sí, por
ser datos ficticios y necesarios para auditar. \estado{verificado}
